\section{Device time, end-to-end time, and hardware-resource accounting}
\label{sec:resources}

\subsection{Timing boundaries actually measured}

The recorded comparison uses, per backend: D-Wave \texttt{qpu\_access\_time} (programming, anneal, readout), excluding network latency, queueing, and one-time embedding; classical-sampler host-wrapper time after excluding JIT compilation warm-up; and FPGA on-device cycle count divided by its $120\,\mathrm{MHz}$ clock. \texttt{src/encoder.py} (line 557) additionally forms \texttt{total\_sampling\_time\_s = sample\_time\_s + cem\_time}, i.e.\ the device/sampler time plus the host-side CEM computation -- a legitimate device-plus-thermometry-overhead number, but still not host-device transfer, embedding/unembedding construction, chain decoding, or the network round-trip. A separate \texttt{wall\_time\_s} covering an entire run is recorded in some pipelines (\texttt{scripts/hparam/hparam\_optuna.py}, \texttt{scripts/ite/ite\_run.py}, \texttt{scripts/j1j2/j1j2\_convergence\_median.py}) but is not decomposed into the components above and is not the metric plotted in the $\mathrm{TTE}_{99}$ figures.

\missing{\textbf{Two clearly labelled timing series.} (1) device/kernel time, approximately what is already plotted; (2) end-to-end wall time including data preparation, host--device transfer, embedding/unembedding, chain decoding, the CEM computation already tracked in \texttt{total\_sampling\_time\_s}, and network round-trips. A third, practically motivated number -- amortized end-to-end time when an embedding is cached and reused across many training jobs -- is also absent. None of these decompositions exist in the current timing code (Sec.~\ref{sec:code-gaps}, item T5), so Fig.~\ref{fig:e2e-timing} [PLACEHOLDER] requires new instrumentation, not a re-plot.}

Stochastic reconfiguration itself is excluded from the sampler-vs-sampler comparison because it is nominally common across backends; that isolates sampler performance but can hide a regime where the SR solve dominates total training time regardless of sampler speed. We do not currently report total training time alongside the sampler-only comparison; doing so requires no new experiments, only re-aggregating already-logged per-iteration timings.

\subsection{Quantum-annealer hardware-resource accounting}
\label{sec:hw-resources}

The reported near-flat scaling of QPU access time with $N$ (fitted exponents $\propto N^{0.16}$ and $\propto N^{0.05}$ over only four sizes, $N\in\{8,16,32,64\}$) reflects a fixed $20\,\mu\mathrm s$ anneal time and a fixed read count per call; it says nothing by itself about how the physical resources consumed by minor embedding grow with $N$, since a dense $K_{N,N}$ RBM requires longer coupler chains as $N$ grows. \texttt{src/sampler.py} (the D-Wave sampling paths, e.g.\ around line 1478 and 1588) reads only \texttt{info["timing"]["qpu\_access\_time"]} from the returned sample set; grepping the repository for chain-break fraction, physical-qubit count, or embedding-success telemetry (\texttt{chain\_break\_fraction}, \texttt{broken\_chain}, \texttt{num\_qubits}, \texttt{embedding\_success}) returns no hits. The existing embedding-algorithm scripts (\texttt{scripts/exper/embedding\_algo\_comparison.py}, \texttt{parallel\_embedding\_experiment.py}, \texttt{parallel\_embedding\_bench.py}) compare how the choice of embedding algorithm affects training outcome and QPU time, not the qubit-count/chain-length/rescaling telemetry needed here, and no programming/anneal/readout time breakdown is recorded anywhere.

Physical qubit count and chain length, at least, do not require new hardware access: the archived result files log an \texttt{embedding\_info} field (present for $N=8,32,64$, both devices; absent for $N=16$) that already gives exactly this, unused until now. The $N=16$ gap is closeable offline: both devices' live hardware-graph snapshots are cached on disk (\texttt{embeddings/\_hwgraph\_\{Advantage\_system6,Advantage2\_system1\}\_live.json}, the same topology used to generate the other sizes' embeddings), so re-running \texttt{minorminer.busclique} locally against that cached graph -- no QPU or network access -- gives a directly comparable number. Table~\ref{tab:resources} reports the complete series.

\begin{table}[h]
\caption{Physical qubits and chain length for the dense $K_{N,N}$ biclique embedding, both devices. $N=16$ recomputed offline from the cached live hardware-graph snapshot (\texttt{minorminer.busclique}, no QPU access); all other cells read from archived \texttt{embedding\_info}.}
\label{tab:resources}
\begin{ruledtabular}
\begin{tabular}{lrrrrr}
& \multicolumn{2}{c}{Pegasus} & \multicolumn{2}{c}{Zephyr} \\
$N$ & qubits & mean chain & qubits & mean chain \\
\hline
8 & 22 & 1.38 & 16 & 1.00 \\
16 & 64 & 2.00 & 48 & 1.50 \\
32 & 212 & 3.31 & 192 & 3.00 \\
64 & 768 & 6.00 & 956 & 7.47 \\
\end{tabular}
\end{ruledtabular}
\end{table}

Zephyr is cheaper than Pegasus at every size up to $N=32$, then crosses over: at $N=64$ -- the size at which the corrected crossover claim of Sec.~\ref{sec:limitations} is made -- Zephyr requires both more physical qubits (956 vs.\ 768) and longer chains (mean $7.47$ vs.\ $6.00$) than Pegasus, a concrete hardware-resource correlate of Zephyr's consistently worse showing there.

\missing{\textbf{Remaining resource fields.} Chain-break fraction, coefficient-rescaling range, embedding success rate, number of simultaneously available embeddings, and separate programming/anneal/readout time components are still not captured anywhere (Sec.~\ref{sec:code-gaps}, item T5) -- these require reading additional fields from the D-Wave SDK's \texttt{SampleSet.info} that the current code discards, which is only obtainable from a live QPU call, not a re-analysis. A more honest timing model, $t(N)=t_0+aN^p$ with $t_0$ a fixed call overhead, should be fit once these remaining resource series exist; with only four $N$ values any such fit remains descriptive.}

\subsection{FPGA sampler: documentation and fidelity}
\label{sec:fpga}

The FPGA backend (\texttt{FPGASampler} in \texttt{src/sampler.py}, lines 581--892) documents its I/O protocol and configuration (\texttt{num\_rep=1024}, \texttt{num\_steps}, \texttt{num\_sweeps}, \texttt{start\_temp}/\texttt{stop\_temp}, geometric schedule) but not its transition rule or stationary distribution; the update algorithm itself lives in the external, non-vendored \texttt{VeloxQFPGA} Julia package, called from \texttt{scripts/fpga/fpga\_sa\_server.jl} rather than reproduced in this repository. How the 1024 replicas become the $\sim\!200$ samples used by SR \emph{is} documented: lines $\approx$858--865 draw a uniform, without-replacement subsample (\texttt{\_subsample\_rng.choice(..., replace=False)}), with an accompanying comment (on the analogous \texttt{VeloxQStandardSASampler}, lines $\approx$891--909) explaining that uniform subsampling, rather than taking the lowest-energy replicas, is required to avoid biasing $\avg{\Eloc}$. No bitstream, board model, or toolchain version is pinned in code (an optional \texttt{FPGA\_BITSTREAM} environment variable defaults to blank), sample independence across replicas is not tested, and no $\DTV$-style fidelity test analogous to Fig.~\ref{fig:classical} exists for the FPGA (\texttt{scripts/dtv/} contains no FPGA-referencing file).

\missing{\textbf{Complete FPGA algorithmic specification and fidelity test.} (i) the transition/update equations actually implemented in \texttt{VeloxQFPGA}, vendored or cited to a pinned version; (ii) whether the resulting chain has a known stationary distribution and how its effective temperature is controlled; (iii) FPGA model, toolchain, and bitstream version, pinned; (iv) a $\DTV$/$\DKL$ fidelity test at small $N$ analogous to Fig.~\ref{fig:classical}, so the FPGA enters the same non-Gibbs-distortion accounting as the QPU and classical samplers rather than being assessed by energy alone.}

\subsection{Energy accounting}

GPU energy is measured, not modeled: \texttt{src/energy.py}'s \texttt{GPUEnergyMeter} polls \texttt{nvidia-smi --query-gpu=power.draw} on a background thread (default $0.5\,\mathrm s$ interval) and trapezoidally integrates power over explicitly marked active windows that isolate solver/sampling time from SR/CG compute; it does not subtract an idle-power baseline. FPGA energy is not measured at all: \texttt{scripts/viz/plot\_ite.py} (lines 82--86) and \texttt{scripts/viz/paper\_figures.py} (line 1348) both hardcode \texttt{ASSUMED\_POWER\_W\{"fpga/fpga": 45.0\}} and multiply by device time; QPU energy is correctly omitted rather than estimated. The current energy comparison therefore mixes one genuinely measured series (GPU) with one modeled constant (FPGA) and one omission (QPU); \missing{it should either move to supplementary material labelled exploratory, or be rebuilt with measured FPGA board power, idle-baseline subtraction for the GPU meter, and a stated common system boundary for both.}
